\chapter{Experimental Evaluation}
\label{experiment}


This chapter evaluates the SoH estimation framework described in Chapter~\ref{methodology}. Section~\ref{sec:exp_setup} details the experimental setup: hardware and software environments, dataset statistics, evaluation metrics, and the BiLSTM baseline configuration. Sections~\ref{sec:exp_hp}--\ref{sec:exp_generalization} present results for each research question in turn: hyperparameter optimization, accuracy comparison against the state of the art, deployment performance across window lengths, and generalization to unseen window lengths. Section~\ref{sec:exp_discussion} discusses the findings and their limitations.


\section{Experimental Setup}
\label{sec:exp_setup}

\subsection{Hardware and Software Environments}

Experiments span two computing environments.

\textbf{Training environment.} Model training and hyperparameter search run on a laptop with an Intel Core i5-11320H CPU (3.20\,GHz) and 16\,GB of RAM. All training is CPU-only. The software stack is Python-based, with LightGBM 4.5.0 for model training and Optuna 4.2.0 for hyperparameter optimization~\autocite{ke2017,akiba2019}. The core packages of the training environment appear in Appendix~\ref{app:dependencies}.

\textbf{Deployment and measurement environment.} The deployment target is the SPC58EC80E5 automotive microcontroller (dual-core Power Architecture e200z420n3 at 180\,MHz, 4\,MB Flash, 384\,KB SRAM)~\autocite{spc58ec80e5_ds}. Code is compiled with the SPC5Studio IDE's integrated GCC Power Architecture cross-compiler at optimization level \texttt{-O2}. Program execution and debugging use UDE Starterkit (PLS).

\textbf{Memory measurement.} Flash and RAM usage are derived from static analysis of the compiled object files: Flash counts the model parameter arrays only (firmware code and HAL are excluded), while RAM counts statically allocated program data. Runtime stack usage is not included, but is bounded and negligible --- tree traversal is implemented iteratively rather than recursively, so the inference call stack is only a few dozen bytes deep. The same measurement methodology is applied to the BiLSTM baseline, keeping the comparison consistent.

Cycle measurements, as described in Section~\ref{sec:ch3_deployment}, use randomly generated window data; accuracy metrics (MAE and MSE) are evaluated in the Python environment on the held-out test set, since the MCU inference output was verified to match the Python output bit-for-bit.

\subsection{Datasets}

Two public NASA lithium-ion battery datasets are used, as introduced in Chapter~\ref{methodology} (Table~\ref{tab:dataset_source}). Table~\ref{tab:dataset_stats} provides statistical details.

\begin{table}[H]
    \centering
    \setcellgapes{3pt}
    \makegapedcells
    \begin{tabular}{c c c}
    \hline
    & NBD & NRD \\
    \hline
    Battery type & 18650 Li-ion & 18650 LCO \\
    Cells & 19 & 28 \\
    Total discharge cycles & 164 & 258 \\
    Sampling rate & 0.09\,Hz & 0.33\,Hz \\
    Label range & 1.29--1.86\,Ah (capacity) & 41.6--100\,\% (SoH) \\
    Train/Val/Test split & \multicolumn{2}{c}{64\% / 16\% / 20\% (cycle-level)} \\
    \hline
    \end{tabular}
    \caption{Dataset statistics. NBD = NASA Battery Dataset; NRD = NASA Randomized Battery Usage Dataset.}
    \label{tab:dataset_stats}
\end{table}

\textbf{Label units.} NBD labels are discharge capacities in ampere-hours (Ah); NRD labels are SoH percentages. In the NBD dataset, the SoH of a cell is expressed through its measured discharge capacity, following the convention of the dataset itself.

\textbf{NBD cell selection.} Following the BiLSTM baseline~\autocite{Sammartino2021}, NBD experiments are restricted to 19 selected cells (battery numbers 5, 6, 7, 18, 25--32, 46--48, 53--56), which excludes cells with abnormal degradation profiles from the original dataset.

\textbf{NRD data collection protocol.} The NRD dataset~\autocite{nasa_data_randomized} was collected by NASA under a protocol that alternates between constant-current characterization discharges (1\,A down to 3.2\,V) and blocks of 50 aging cycles in which the load current is randomly drawn from $\{0.5, 1, 1.5, 2, 2.5, 3, 3.5, 4\}$\,A every 5 minutes. This design produces irregular, multi-rate discharge curves that more faithfully represent the stochastic nature of real-world driving than the constant-current profiles of the NBD dataset.

\textbf{Data split.} The split strategy follows the procedure described in Chapter~\ref{methodology} (Section~\ref{sec:ch3_data}): cycle-level partitioning with 64/16/20 train/validation/test proportions, with cross-cell isolation enforced for NRD.

\subsection{Evaluation Metrics}

The metrics defined in Chapter~\ref{methodology} (Section~\ref{sec:ch3_verification}) are briefly restated here:

\begin{itemize}
    \item \textbf{MAE} (Mean Absolute Error): $\frac{1}{N}\sum |\hat{y}_i - y_i|$. Primary metric.
    \item \textbf{MSE} (Mean Squared Error): $\frac{1}{N}\sum (\hat{y}_i - y_i)^2$. Secondary metric, penalizes large errors.
    \item \textbf{FW} (First Window): MAE computed on the first window of each test discharge cycle only, following the evaluation protocol of Sammartino et al.~\autocite{Sammartino2021}.
    \item \textbf{Full}: MAE on all test windows.
    \item \textbf{Inference cycles}: CPU clock cycles for one complete inference (feature extraction + model traversal), measured with the MCU's system timer (STMD3) as described in Section~\ref{sec:ch3_deployment}.
    \item \textbf{Flash / RAM}: Static memory usage derived from the compiled object files, reported in absolute bytes and as a percentage of the 4\,MB Flash and 384\,KB SRAM.
\end{itemize}

\subsection{BiLSTM Baseline}

To ensure a fair comparison, the BiLSTM model of Sammartino et al.~\autocite{Sammartino2021} is reproduced. This baseline is one of the few published works covering the full training-to-deployment pipeline for SoH estimation on a 120\,MHz SPC584B MCU, and it uses the same NBD dataset.

The BiLSTM was retrained using the same data split as the LightGBM model. Configuration details:

\begin{itemize}
    \item Architecture: a single bidirectional LSTM layer with hidden size 14
    \item Training: 45 epochs, learning rate 0.01, batch size 16
    \item Input window: $wlen = 20$ (fixed; original paper setting)
    \item MCU deployment: Keras model converted to C via SPC5Studio AI
\end{itemize}

For the multi-window-length deployment experiment (Section~\ref{sec:exp_deployment}), a separate BiLSTM model is trained for each $wlen \in \{20, 30, 40, 50, 60\}$, because the SPC5Studio AI conversion fixes the input tensor shape. For the generalization experiment (Section~\ref{sec:exp_generalization}), unseen window lengths are handled by trimming the oldest time steps for inputs longer than the training length, and zero-padding for shorter sequences. In contrast, the LightGBM pipeline accommodates arbitrary window lengths without modification, because the feature extraction step always produces the same fixed-size vector.


\section{Hyperparameter Optimization Results}
\label{sec:exp_hp}

After 1{,}000 Optuna trials, the best hyperparameter configurations for each dataset are shown in Table~\ref{tab:best_hp}.

\begin{table}[H]
    \centering
    \setcellgapes{3pt}
    \makegapedcells
    \begin{tabular}{c c c c c c c c c}
    \hline
    Dataset & LR & \makecell{Max \\ Depth} & Leaves & Estimators & $\alpha$ & $\beta$ & MCS & CST \\
    \hline
    NBD & 0.147 & 8 & 25 & 384 & 6.5e-2 & 6.7e-5 & 6 & 0.46 \\
    NRD & 0.085 & 8 & 18 & 443 & 1.9e-6 & 3.6e-4 & 17 & 0.61 \\
    \hline
    \end{tabular}
    \caption{Best hyperparameter configurations found by Bayesian search}
    \label{tab:best_hp}
\end{table}

Several observations stand out. First, the optimal learning rate for NRD (0.085) is substantially lower than for NBD (0.147), while the ensemble is larger (443 trees vs.\ 384). NRD's randomized charge/discharge profiles and multi-cell composition produce more complex feature distributions, requiring smaller gradient steps and more base learners.

Second, NRD's leaf count (18) is lower than NBD's (25), and the search converged to a more conservative configuration --- shallower trees and a five-times-stronger L2 penalty, although the L1 penalty is orders of magnitude weaker --- to prevent overfitting on the noisier data.

Third, \texttt{max\_depth} hit the upper bound of 8 on both datasets. Without this deployment-motivated constraint, the search would likely have explored deeper trees. As argued in Chapter~\ref{methodology}, the cap ensures the resulting model fits comfortably within the MCU's Flash budget, and the accuracy results in the following sections confirm that the constraint does not materially harm estimation quality.


\section{Accuracy Comparison with BiLSTM}
\label{sec:exp_sota}

Table~\ref{tab:sota} compares LightGBM against BiLSTM at $wlen = 20$, the window length for which the baseline was originally trained. Alongside accuracy (MAE FW and Full), the table reports inference latency and estimated energy per inference.

\begin{table}[H]
    \centering
    \setcellgapes{3pt}
    \makegapedcells
    \begin{tabular}{c c c c c}
    \hline
    Model & \makecell{MAE (FW) \\ \relax[Ah]} & \makecell{MAE (Full) \\ \relax[Ah]} & \makecell{Latency \\ \relax[ms]} & \makecell{Estimated energy \\ \relax[$\mu$J]} \\
    \hline
    BiLSTM~\autocite{Sammartino2021} & 0.0202 & 0.0194 & 4.656 & 279.36 \\
    LightGBM (Ours) & 0.0167 & 0.0188 & 0.403 & 24.18 \\
    \hline
    \end{tabular}
    \caption{State-of-the-art comparison on the NBD dataset at $wlen = 20$. MAE is reported in Ah; Table~\ref{tab:deployment} reports the same comparison on the normalized SoH scale.}
    \label{tab:sota}
\end{table}

On accuracy, LightGBM reduces FW MAE by approximately 17\% (0.0202 $\rightarrow$ 0.0167) and Full MAE by about 3\% (0.0194 $\rightarrow$ 0.0188).

On efficiency, the gap is larger. LightGBM inference takes 0.403\,ms versus 4.656\,ms for BiLSTM --- an 11.5$\times$ speedup. The estimated energy per inference drops from 279.36\,$\mu$J to 24.18\,$\mu$J. SoH estimation in a real BMS runs once per window (not once per sample), so the absolute energy cost is negligible in either case. But the latency reduction frees a meaningful amount of CPU time in the BMS main loop, which must also handle cell balancing, thermal monitoring, communication, and safety checks within hard real-time deadlines.

The energy figures are derived as $E = P \times t_{inf}$, with the core power $P$ estimated from the device datasheet (DS11620): the typical core supply voltage $V_{DD\_LV}$ is 1.2\,V and the maximum main-core dynamic current $I_{DD\_MAIN\_CORE\_AC}$ is 50\,mA, giving $P = 1.2\,\text{V} \times 50\,\text{mA} = 60$\,mW~\autocite{spc58ec80e5_ds}. Because the current figure is a maximum, the energy values are a conservative upper bound. The same convention is applied to the BiLSTM baseline for comparability.


\section[Deployment Performance Across Window Lengths]{Deployment Performance Across Window\\Lengths}
\label{sec:exp_deployment}

Table~\ref{tab:deployment} is the central result of this chapter. It reports MAE, MSE, inference cycles, RAM usage, and Flash usage for both models across all five training window lengths on both datasets.

\subsection{Accuracy}

On the NBD dataset, the two methods are close at short window lengths ($wlen \leq 30$), with LightGBM holding a small edge (MAE difference below 0.04 percentage points). From $wlen = 40$ onward, BiLSTM accuracy degrades markedly: MAE jumps from 0.87\% at $wlen = 30$ to 2.16\% at $wlen = 40$, a 2.5$\times$ increase, and remains elevated at longer windows. LightGBM stays stable within a narrow band of 0.74--1.01\% across all window lengths.

\begin{table}[H]
    \centering
    \setcellgapes{3pt}
    \makegapedcells
    \begin{tabular}{c c c c c c c c}
    \hline
    Dataset & WL & Model & \makecell{MAE \\ \relax[\%]} & \makecell{MSE \\ \relax[\%$^2$]} & \makecell{Cycles \\ \relax[kcycles]} & \makecell{RAM \\ \relax[\%]} & \makecell{Flash \\ \relax[\%]} \\
    \hline
    \multirow{10}{*}{NBD}
    & 20 & L & 1.013 & 2.321 & 72.6 & 0.49 & 2.71 \\
    & 20 & B & 1.045 & 2.031 & 838 & 0.13 & 0.21 \\
    & 30 & L & 0.840 & 1.741 & 74.7 & 0.49 & 2.71 \\
    & 30 & B & 0.873 & 1.451 & 1250 & 0.13 & 0.21 \\
    & 40 & L & 0.738 & 1.161 & 77.1 & 0.49 & 2.71 \\
    & 40 & B & 2.160 & 7.544 & 1670 & 0.13 & 0.21 \\
    & 50 & L & 0.921 & 2.321 & 79.0 & 0.49 & 2.71 \\
    & 50 & B & 1.551 & 4.062 & 2080 & 0.13 & 0.21 \\
    & 60 & L & 0.932 & 1.741 & 81.3 & 0.49 & 2.71 \\
    & 60 & B & 1.514 & 3.772 & 2500 & 0.13 & 0.21 \\
    \hline
    \multirow{10}{*}{NRD}
    & 20 & L & 2.256 & 12.622 & 73.3 & 0.49 & 2.23 \\
    & 20 & B & 2.908 & 14.897 & 838 & 0.13 & 0.21 \\
    & 30 & L & 2.208 & 11.788 & 75.4 & 0.49 & 2.23 \\
    & 30 & B & 4.127 & 27.588 & 1250 & 0.13 & 0.21 \\
    & 40 & L & 2.183 & 10.498 & 77.8 & 0.49 & 2.23 \\
    & 40 & B & 4.002 & 28.180 & 1670 & 0.13 & 0.21 \\
    & 50 & L & 1.881 & 7.137 & 79.6 & 0.49 & 2.23 \\
    & 50 & B & 4.751 & 40.870 & 2080 & 0.13 & 0.21 \\
    & 60 & L & 1.751 & 7.357 & 82.0 & 0.49 & 2.23 \\
    & 60 & B & 4.368 & 30.323 & 2500 & 0.13 & 0.21 \\
    \hline
    \end{tabular}
    \caption{Deployment results. L = LightGBM (this work), B = BiLSTM, WL = Window Length. MAE and MSE are computed on the SoH scale: NBD capacities are divided by the maximum capacity (1.8565\,Ah) before evaluation, so both datasets report errors in percent and percent-squared; both are evaluated over the full set of test windows (Full protocol). Cycles are in thousands (kcycles). RAM and Flash are percentages of the SPC58EC80E5's 384\,KB SRAM and 4\,MB Flash.}
    \label{tab:deployment}
\end{table}

On the more complex NRD dataset, LightGBM's advantage is both larger and consistent across window lengths. At $wlen = 20$, LightGBM already leads ($\text{MAE} = 2.26$ vs.\ 2.91), and the gap widens steadily with window length. At $wlen = 60$, LightGBM achieves a MAE roughly 2.5$\times$ lower than BiLSTM. Moreover, LightGBM's MAE actually \emph{decreases} with longer windows (from 2.26 to 1.75), while BiLSTM's \emph{increases} (from 2.91 to 4.37). This seemingly counterintuitive behavior has a straightforward explanation. Longer windows give the feature extraction step more samples to compute its statistical aggregates (means, min, max), yielding more stable feature estimates and therefore better inputs to the model. BiLSTM, processing raw time series, cannot exploit this ``statistical convergence'' effect. Instead, longer sequences aggravate the difficulty of capturing long-range temporal dependencies.



\subsection{Inference Efficiency}

The two methods show qualitatively different scaling behavior.

LightGBM's inference cycles grow from approximately 73k at $wlen = 20$ to 82k at $wlen = 60$, an increase of about 12\%. Nearly all of this increase comes from feature extraction --- computing means and extrema over a longer buffer takes more loop iterations. The inference cost of the model component (tree traversal) is constant with window length, since the model always receives a 10-element vector.

BiLSTM's inference cycles grow from 838k at $wlen = 20$ to 2.50M at $wlen = 60$, an increase of approximately 3$\times$. This is a structural consequence of recurrent architectures: every time step requires matrix multiplications and activation functions, and the cost scales linearly with the number of steps.

In terms of speedup, LightGBM is 11.5$\times$ faster at $wlen = 20$ and 30.7$\times$ faster at $wlen = 60$. The advantage grows with window length and, by extension, with sampling rate. In a BMS operating at higher sampling frequencies, the gap would be even wider.

\subsection{Memory Footprint}

LightGBM uses approximately 2.2--2.7\% of Flash ($\approx$ 91--111\,KB) and 0.49\% of RAM ($\approx$ 1.88\,KB). BiLSTM uses 0.21\% Flash and 0.13\% RAM.

LightGBM is heavier on memory, primarily because the tree ensemble's four arrays occupy Flash. The parameter file is composed of four arrays: \texttt{uint16\_t ROOTS[N\_TREES]} (tree root indices), \texttt{uint8\_t FEATURES[N\_NODES]} (feature index or leaf marker per node), \texttt{float ALPHAS[N\_NODES]} (threshold or leaf value per node), and \texttt{uint8\_t CHILDREN\_RIGHT[N\_NODES]} (right child index per node). Each tree node therefore requires 6 bytes (48 bits) across the three per-node arrays; ROOTS adds a negligible 2 bytes per tree.

BiLSTM's lighter memory profile reflects the optimization of SPC5Studio AI's neural network runtime, which uses compact weight representations. However, this efficiency comes at the cost of limited model-type support: SPC5Studio AI only handles specific neural network architectures, whereas the EDEN-compiled LightGBM code is self-contained and portable.

In absolute terms, both models are well within the SPC58's memory budget: the compiled LightGBM model occupies 111\,KB (NBD) and 91.34\,KB (NRD) of the 4\,MB Flash. Allocating 2--3\% of Flash to the model in exchange for an 11--31$\times$ inference speedup and better or equivalent accuracy is a clearly favorable trade-off for this hardware platform.


\section{Generalization to Unseen Window Lengths}
\label{sec:exp_generalization}

The experiments so far evaluated performance on window lengths seen during training. A more stringent test is: can the model handle window lengths it was never trained on? This simulates real-world contingencies --- a BMS firmware update changes the sampling rate, or an unusually short trip produces a window far from the nominal configuration.

\begin{table}[H]
    \centering
    \renewcommand{\arraystretch}{1.3}
    \begin{minipage}[t]{0.45\linewidth}
    \centering
    \begin{tabular}{l c c c c}
    \hline
    & 10 & 35 & 55 & 75 \\
    \hline
    L & \hc{hg10}{0.041} & \hc{hg20}{0.023} & \hc{hg20}{0.025} & \hc{hy08}{0.095} \\
    B$_{20}$ & \hc{hy16}{0.106} & \hc{hr50}{0.491} & \hc{hr50}{0.492} & \hc{hy24}{0.274} \\
    B$_{30}$ & \hc{hy12}{0.059} & \hc{hg15}{0.042} & \hc{hr55}{0.504} & \hc{hr55}{0.497} \\
    B$_{40}$ & \hc{hy16}{0.102} & \hc{hy08}{0.073} & \hc{hr60}{0.668} & \hc{hr60}{0.660} \\
    B$_{50}$ & \hc{hy16}{0.097} & \hc{hy08}{0.068} & \hc{hg10}{0.033} & \hc{hr50}{0.643} \\
    B$_{60}$ & \hc{hg10}{0.041} & \hc{hg20}{0.023} & \hc{hg20}{0.025} & \hc{hy08}{0.095} \\
    \hline
    \end{tabular}\par\vspace{2pt}
    \textbf{(a) NBD}
    \end{minipage}\hfill
    \begin{minipage}[t]{0.45\linewidth}
    \centering
    \begin{tabular}{l c c c c}
    \hline
    & 10 & 35 & 55 & 75 \\
    \hline
    L & \hc{hg15}{3.7} & \hc{hg25}{2.3} & \hc{hg30}{1.7} & \hc{hg25}{2.3} \\
    B$_{20}$ & \hc{hy16}{5.3} & \hc{hr70}{15.3} & \hc{hy32}{7.5} & \hc{hy36}{10.8} \\
    B$_{30}$ & \hc{hy16}{5.5} & \hc{hr75}{19.8} & \hc{hr75}{19.9} & \hc{hy32}{12.6} \\
    B$_{40}$ & \hc{hy24}{6.8} & \hc{hg15}{4.2} & \hc{hr70}{24.0} & \hc{hr75}{24.5} \\
    B$_{50}$ & \hc{hy32}{7.5} & \hc{hy20}{4.8} & \hc{hr70}{21.7} & \hc{hr75}{22.2} \\
    B$_{60}$ & \hc{hy36}{8.9} & \hc{hy20}{4.8} & \hc{hg15}{4.2} & \hc{hr80}{28.2} \\
    \hline
    \end{tabular}\par\vspace{2pt}
    \textbf{(b) NRD}
    \end{minipage}
    \caption{MAE of the LightGBM model and the five BiLSTM variants on unseen window lengths, for both datasets. L = LightGBM; B = BiLSTM (subscript: training window length). Cell color encodes the MAE value: green, yellow and red indicate low, medium and high error, respectively; the color scale is defined separately for each dataset. MAE values are in Ah for NBD and in \% for NRD. The LightGBM row stays within the low-error (green to pale yellow) range across all test lengths, while the BiLSTM rows show erratic, length-dependent degradation.}
    \label{tab:generalization_heatmap}
\end{table}

The experiment uses four test lengths not present in the training set --- namely $\{10, 35, 55, 75\}$ --- where 10 is shorter than the minimum training length (20), 75 is longer than the maximum (60), and 35 and 55 fall within the training range but were never explicitly seen. LightGBM requires no modification to handle these lengths, thanks to the fixed-size feature vector produced by feature extraction. The deployed BiLSTM, whose input tensor shape is fixed by the SPC5Studio AI conversion, requires heuristic adaptation instead: inputs longer than the training length are trimmed by discarding the oldest time steps, and shorter sequences are zero-padded, as described in Section~\ref{sec:exp_setup}.

Table~\ref{tab:generalization_heatmap} provides an overview of the generalization results for both datasets; the detailed MAE values are reported in the following subsections.



\subsection{NBD Results}

Table~\ref{tab:generalization_nbd} reports the MAE values for unseen window lengths on the NBD dataset.

\begin{table}[H]
    \centering
    \setcellgapes{3pt}
    \makegapedcells
    \begin{tabular}{c c c c c c}
    \hline
    \multirow{2}{*}{Model} & \multicolumn{4}{c}{Test Window Length} & \multirow{2}{*}{Avg. MAE} \\
    \cline{2-5}
    & 10 & 35 & 55 & 75 & \\
    \hline
    LightGBM (Ours) & 0.041 & 0.023 & 0.025 & 0.095 & 0.046 \\
    BiLSTM$_{20}$ & 0.106 & 0.491 & 0.492 & 0.274 & 0.341 \\
    BiLSTM$_{30}$ & 0.059 & 0.042 & 0.504 & 0.497 & 0.276 \\
    BiLSTM$_{40}$ & 0.102 & 0.073 & 0.668 & 0.660 & 0.376 \\
    BiLSTM$_{50}$ & 0.097 & 0.068 & 0.033 & 0.643 & 0.210 \\
    BiLSTM$_{60}$ & 0.041 & 0.023 & 0.025 & 0.095 & 0.046 \\
    \hline
    \end{tabular}
    \caption{MAE on unseen window lengths  ---  NBD dataset. BiLSTM$_{XX}$ = BiLSTM trained on window length XX. All models evaluated on the same test split. MAE in Ah.}
    \label{tab:generalization_nbd}
\end{table}

LightGBM achieves an average MAE of 0.046 across the four unseen lengths, with reasonably stable performance at most test points, though at $wlen = 75$ the error rises to 0.095 --- about five times the worst seen-length error, but still far below the BiLSTM averages. The five BiLSTM models, each trained on a single length, produce erratic results. BiLSTM$_{40}$ reports MAE 0.668\,Ah at $wlen = 55$, nearly 17$\times$ worse than the 0.040\,Ah (2.16\,\% on the SoH scale) it achieves at its training length. There is no consistent relationship between training window length and generalization performance: models do not necessarily perform best at test lengths near their training configuration.

\subsection{NRD Results}

Table~\ref{tab:generalization_nrd} reports the equivalent results on NRD.

On the more complex NRD data, the contrast is sharper. LightGBM averages MAE 2.50 with small variance across test lengths. The BiLSTM models exhibit extreme fluctuations. The most striking case is BiLSTM$_{60}$: MAE is 4.2 at $wlen = 55$ but jumps to 28.2 at $wlen = 75$ --- a nearly 7$\times$ increase between two test lengths only 20 samples apart. BiLSTM$_{30}$ performs decently at $wlen = 10$ (MAE 5.5) and collapses at $wlen = 55$ (19.9). Across all BiLSTM variants, the average MAE ranges from 9.73 to 14.88 --- approximately 4--6$\times$ worse than LightGBM.

\begin{table}[H]
    \centering
    \setcellgapes{3pt}
    \makegapedcells
    \begin{tabular}{c c c c c c}
    \hline
    \multirow{2}{*}{Model} & \multicolumn{4}{c}{Test Window Length} & \multirow{2}{*}{Avg. MAE} \\
    \cline{2-5}
    & 10 & 35 & 55 & 75 & \\
    \hline
    LightGBM (Ours) & 3.7 & 2.3 & 1.7 & 2.3 & 2.50 \\
    BiLSTM$_{20}$ & 5.3 & 15.3 & 7.5 & 10.8 & 9.73 \\
    BiLSTM$_{30}$ & 5.5 & 19.8 & 19.9 & 12.6 & 14.45 \\
    BiLSTM$_{40}$ & 6.8 & 4.2 & 24.0 & 24.5 & 14.88 \\
    BiLSTM$_{50}$ & 7.5 & 4.8 & 21.7 & 22.2 & 14.05 \\
    BiLSTM$_{60}$ & 8.9 & 4.8 & 4.2 & 28.2 & 11.53 \\
    \hline
    \end{tabular}
    \caption{MAE on unseen window lengths  ---  NRD dataset. MAE in \,\%.}
    \label{tab:generalization_nrd}
\end{table}

\subsection{Why LightGBM Generalizes}

The generalization results follow directly from the design documented in Chapter~\ref{methodology}. LightGBM receives a fixed 10-dimensional feature vector, not a raw time series. Feature extraction computes the same statistics regardless of window length. The multi-length training strategy (Section~\ref{sec:ch3_segmentation}) then teaches the model that the feature distribution may shift slightly with window length, but the SoH label should remain unchanged. The model learns this invariance, and it holds when new lengths appear.

BiLSTM's failure mode is structural. Changing the input sequence length by truncation or padding creates artifacts the model never encountered during training. For BiLSTM$_{40}$, trained on sequences of exactly 40 steps, a 55-step sequence truncated to 40 discards the 15 oldest samples --- if those samples contained the early portion of the voltage drop, the model has lost information it was trained to expect. A 10-step sequence padded to 40 injects 30 artificial zeros, which the model's recurrent dynamics were never calibrated to handle. LightGBM sidesteps this entire class of problems by operating on features rather than sequences.


\section{Discussion and Limitations}
\label{sec:exp_discussion}

\subsection{Summary of Findings}

Three central findings emerge from the experiments.

\textbf{Accuracy.} The LightGBM + multi-length feature extraction scheme matches or exceeds BiLSTM accuracy on both NBD and NRD datasets. The advantage is most pronounced on the more complex, multi-cell NRD dataset and at longer window lengths. This validates the analysis in Chapter~\ref{background}: tree ensembles extracting statistical features from variable-length windows capture degradation information more robustly than recurrent networks processing fixed-length sequences.

\textbf{Efficiency.} LightGBM inference is 11--31$\times$ faster than BiLSTM, and the cost is nearly independent of window length. This property is critical for deployment in BMS hardware with variable sampling rates. BiLSTM's linearly scaling cost poses a real-time risk under longer windows or higher sampling frequencies.

\textbf{Robustness.} A single LightGBM model, trained once, maintains stable accuracy on window lengths it never encountered during training. BiLSTM under the same conditions produces unpredictable and often severely degraded accuracy. In practical terms, this means the LightGBM model can handle unforeseen changes in operating conditions over the vehicle's lifetime without retraining.

\subsection{Limitations}

Several limitations should be considered when interpreting these results.

First, both datasets come from laboratory environments. Charge and discharge cycles are conducted under controlled temperatures with pre-programmed load profiles. Real-world batteries experience wider temperature swings ($-20^\circ$C cold starts to $+50^\circ$C summer heat), irregular load patterns (acceleration, regenerative braking, idling), and mixed aging histories. Validation on data collected from fielded vehicles remains an important direction for future work.

Second, the experiments evaluate only the SoH component of the dual-model architecture described in Chapter~\ref{introduction}. The SoC model and the OTA update mechanism are not experimentally validated here. Preliminary SoC experiments confirmed that SoH estimation errors propagate to SoC predictions, but a complete solution to this coupling problem remains for future work.

Third, MCU performance measurements use randomly generated data and report pure computation time. In a production BMS, data transfer from the analog front-end to the MCU, interrupt handling, and DMA operations add to the end-to-end latency. The cycle counts reported here represent a lower bound on real-world inference time.

Fourth, the 64/16/20 split and NRD cross-cell isolation ensure statistical validity, but the NBD dataset is relatively small. Evaluation on larger, more diverse datasets would strengthen the generalizability of these conclusions.
